proceeding by descending induction, suppose it is known for @@M0@@. After replacing @@M1@@ by @@M2@@, we may suppose that @@M3@@ is @@M4@@-flat for @@M5@@. Since @@M6@@ is flat over @@M7@@, it then follows from EGA III 6.9.8 that @@M8@@ commutes with every base change. Consequently, the representability of @@M9@@ follows from the next lemma (see: J.P. Murre, Representation of unramified functors, Séminaire Bourbaki; May 1965, p.294-11, th.2):

\sgaStableTarget{sga6.e12.lem.4-5}\textbf{Lemma 4.5.} \uline{Let} @@M10@@ \uline{be a proper morphism of finite presentation, and} @@M11@@ \uline{a sheaf of finite presentation on} @@M12@@\uline{. Then the subfunctor} @@M13@@ \uline{of} @@M14@@ \uline{that makes} @@M15@@ \uline{flat over} @@M16@@ \uline{is representable by a surjective monomorphism of finite presentation (which is of course not an immersion in general).}

To represent condition 3) of 4.3, note that, by 1) and 2), @@M17@@ and @@M18@@ commute with base changes @@M19@@, and hence @@M20@@ does also; then apply 4.5.

To represent 4), note that, by 3) and the fact that @@M21@@ and @@M22@@ are flat and commute with base changes, @@M23@@ also commutes; then apply the

\sgaStableTarget{sga6.e12.lem.4-6}\textbf{Lemma 4.6.} \uline{Let} @@M24@@ \uline{be a proper morphism of finite presentation, and} @@M25@@ \uline{a homomorphism of} @@M26@@\uline{-Modules of finite presentation, with} @@M27@@ \uline{flat over} @@M28@@\uline{. Then the subfunctor} @@M29@@ \uline{of} @@M30@@ \uline{that makes} @@M31@@ \uline{an isomorphism is representable by an open subscheme of} @@M32@@\uline{, of finite presentation.}

As usual, we reduce to the case where @@M33@@ is noetherian, using EGA IV 8 and 11.2.6, which relieves us of the need to consider finite presentation. First let us represent the functor that makes @@M34@@ an epimorphism,

